% GENERATED FILE. Edit canonical lessons, then run npm run generate:books.
% canonical-source-hash: fnv1a64:9cc07e3d45436dc0
% canonical-lessons: TE-C519-cikitsa, TE-C519-sastracikitsa, TE-C519-mattumandu, TE-C519-patti, TE-C519-paramarsincu

\chapter{Treatment}
\label{ch:te-treatment}
\hlchaptermodality{519}

\begin{chapteropening}
\textbf{By the end of this chapter:} \emph{I can talk about treatment, surgery and visiting the sick.}

Complete the last lesson of chapter 519: I can talk about treatment, surgery and visiting the sick.
\end{chapteropening}

\section[\texorpdfstring{\te{చికిత్స} (cikitsa)}{చికిత్స (cikitsa)}]{\te{చికిత్స} (cikitsa) — treatment; therapy}

\label{lesson:TE-C519-cikitsa}



\begin{quote}
Before the new one: say the Telugu for to limp, then the Telugu for bleeding.
\end{quote}

\subsection*{You'll want to know: \te{చికిత్స}}
\textbf{\te{చికిత్స}} — \emph{cikitsa} — \textquotedblleft{}treatment; therapy\textquotedblright{}.

\begin{quote}
\te{మంచి} \te{చికిత్స} \te{వల్ల} \te{ఆమె} \te{త్వరగా} \te{కోలుకుంది}.
\end{quote}

\begin{quote}
\emph{mañci cikitsa valla āme tvaragā kōlukundi.}
\end{quote}

\begin{quote}
Thanks to good treatment she recovered quickly.
\end{quote}

\te{చికిత్స} \te{చేయు}, to treat; \te{చికిత్స} \te{పొందు}, to receive treatment.

\begin{cousinweb}
\textbf{\te{చికిత్స}} (\emph{cikitsa}) is borrowed from Sanskrit \emph{cikitsā}, \textquotedblleft{}medical attendance, the practice of medicine\textquotedblright{} (Monier-Williams).
\end{cousinweb}

\begin{grammarlens}[title={What you've built}]
One more word for describing what you lived and felt.
\end{grammarlens}

\subsection*{Your turn}
\begin{itemize}
\raggedright
  \item \emph{Say it:} \emph{cikitsa}
  \item \emph{Say it:} \emph{cikitsa}, once more
  \item \emph{Say it:} \emph{mañci cikitsa valla āme tvaragā kōlukundi.} — Thanks to good treatment she recovered quickly.
  \item \emph{Recall:} say the Telugu for to limp, then the Telugu for bleeding, then say \emph{cikitsa} again
\end{itemize}

\begin{tcolorbox}[breakable,colback=teal!4,colframe=teal!35!black,title={Before you move on}]
What does \te{చికిత్స} mean? (\textquotedblleft{}Treatment; therapy\textquotedblright{}.) Say it once more.
\end{tcolorbox}
\section[\texorpdfstring{\te{శస్త్రచికిత్స} (śastracikitsa)}{శస్త్రచికిత్స (śastracikitsa)}]{\te{శస్త్రచికిత్స} (śastracikitsa) — surgery; an operation}

\label{lesson:TE-C519-sastracikitsa}



\begin{quote}
Before the new one: say the Telugu for bleeding, then the Telugu for treatment.
\end{quote}

\subsection*{You'll want to know: \te{శస్త్రచికిత్స}}
\textbf{\te{శస్త్రచికిత్స}} — \emph{śastracikitsa} — \textquotedblleft{}surgery; an operation\textquotedblright{}.

\begin{quote}
\te{తాతయ్యకు} \te{మోకాలికి} \te{శస్త్రచికిత్స} \te{జరిగింది}.
\end{quote}

\begin{quote}
\emph{tātayyaku mōkāliki śastracikitsa jarigindi.}
\end{quote}

\begin{quote}
Grandpa had surgery on his knee.
\end{quote}

\te{శస్త్రం}, blade, + \te{చికిత్స}: the formal word for \te{ఆపరేషను}.

\begin{cousinweb}
\textbf{\te{శస్త్రచికిత్స}} (\emph{śastracikitsa}) is a learned compound built from Sanskrit \emph{śastra}, \textquotedblleft{}an instrument for cutting\textquotedblright{} (CDIAL 12367), and \emph{cikitsā}, \textquotedblleft{}medical treatment\textquotedblright{} (Monier-Williams).
\end{cousinweb}

\begin{grammarlens}[title={What you've built}]
One more word for describing what you lived and felt.
\end{grammarlens}

\subsection*{Your turn}
\begin{itemize}
\raggedright
  \item \emph{Say it:} \emph{śastracikitsa}
  \item \emph{Say it:} \emph{śastracikitsa}, once more
  \item \emph{Say it:} \emph{tātayyaku mōkāliki śastracikitsa jarigindi.} — Grandpa had surgery on his knee.
  \item \emph{Recall:} say the Telugu for bleeding, then the Telugu for treatment, then say \emph{śastracikitsa} again
\end{itemize}

\begin{tcolorbox}[breakable,colback=teal!4,colframe=teal!35!black,title={Before you move on}]
What does \te{శస్త్రచికిత్స} mean? (\textquotedblleft{}Surgery; an operation\textquotedblright{}.) Say it once more.
\end{tcolorbox}
\section[\texorpdfstring{\te{మత్తుమందు} (mattumandu)}{మత్తుమందు (mattumandu)}]{\te{మత్తుమందు} (mattumandu) — an anaesthetic}

\label{lesson:TE-C519-mattumandu}



\begin{quote}
Before the new one: say the Telugu for treatment, then the Telugu for surgery.
\end{quote}

\subsection*{You'll want to know: \te{మత్తుమందు}}
\textbf{\te{మత్తుమందు}} — \emph{mattumandu} — \textquotedblleft{}an anaesthetic\textquotedblright{}.

\begin{quote}
\te{ఆపరేషనుకు} \te{ముందు} \te{డాక్టరు} \te{మత్తుమందు} \te{ఇచ్చారు}.
\end{quote}

\begin{quote}
\emph{āparēṣanuku mundu ḍākṭaru mattumandu iccāru.}
\end{quote}

\begin{quote}
Before the operation the doctor gave an anaesthetic.
\end{quote}

\te{మత్తు}, numbness or a stupor, + \te{మందు}, medicine.

\begin{cousinweb}
\textbf{\te{మత్తుమందు}} (\emph{mattumandu}) joins \te{మత్తు}, \textquotedblleft{}a stupor\textquotedblright{}, from Sanskrit \emph{matta}, \textquotedblleft{}intoxicated\textquotedblright{} (CDIAL 9750), to native \emph{mandu}, \textquotedblleft{}medicine\textquotedblright{} (DEDR 4719).
\end{cousinweb}

\begin{grammarlens}[title={What you've built}]
One more word for describing what you lived and felt.
\end{grammarlens}

\subsection*{Your turn}
\begin{itemize}
\raggedright
  \item \emph{Say it:} \emph{mattumandu}
  \item \emph{Say it:} \emph{mattumandu}, once more
  \item \emph{Say it:} \emph{āparēṣanuku mundu ḍākṭaru mattumandu iccāru.} — Before the operation the doctor gave an anaesthetic.
  \item \emph{Recall:} say the Telugu for treatment, then the Telugu for surgery, then say \emph{mattumandu} again
\end{itemize}

\begin{tcolorbox}[breakable,colback=teal!4,colframe=teal!35!black,title={Before you move on}]
What does \te{మత్తుమందు} mean? (\textquotedblleft{}An anaesthetic\textquotedblright{}.) Say it once more.
\end{tcolorbox}
\section[\texorpdfstring{\te{పట్టీ} (paṭṭī)}{పట్టీ (paṭṭī)}]{\te{పట్టీ} (paṭṭī) — a bandage; a strip}

\label{lesson:TE-C519-patti}



\begin{quote}
Before the new one: say the Telugu for surgery, then the Telugu for an anaesthetic.
\end{quote}

\subsection*{You'll want to know: \te{పట్టీ}}
\textbf{\te{పట్టీ}} — \emph{paṭṭī} — \textquotedblleft{}a bandage; a strip\textquotedblright{}.

\begin{quote}
\te{నర్సు} \te{గాయానికి} \te{పట్టీ} \te{కట్టింది}.
\end{quote}

\begin{quote}
\emph{narsu gāyāniki paṭṭī kaṭṭindi.}
\end{quote}

\begin{quote}
The nurse bandaged the wound.
\end{quote}

\te{పట్టీ} \te{కట్టు}, to bandage. Also a strip of cloth or a band.

\begin{cousinweb}
\textbf{\te{పట్టీ}} (\emph{paṭṭī}) is usually traced to Sanskrit \emph{paṭṭa}, \textquotedblleft{}cloth; a bandage, a fillet\textquotedblright{} (CDIAL 7700).
\end{cousinweb}

\begin{grammarlens}[title={What you've built}]
One more word for describing what you lived and felt.
\end{grammarlens}

\subsection*{Your turn}
\begin{itemize}
\raggedright
  \item \emph{Say it:} \emph{paṭṭī}
  \item \emph{Say it:} \emph{paṭṭī}, once more
  \item \emph{Say it:} \emph{narsu gāyāniki paṭṭī kaṭṭindi.} — The nurse bandaged the wound.
  \item \emph{Recall:} say the Telugu for surgery, then the Telugu for an anaesthetic, then say \emph{paṭṭī} again
\end{itemize}

\begin{tcolorbox}[breakable,colback=teal!4,colframe=teal!35!black,title={Before you move on}]
What does \te{పట్టీ} mean? (\textquotedblleft{}A bandage; a strip\textquotedblright{}.) Say it once more.
\end{tcolorbox}
\section[\texorpdfstring{\te{పరామర్శించు} (parāmarśiñcu)}{పరామర్శించు (parāmarśiñcu)}]{\te{పరామర్శించు} (parāmarśiñcu) — to visit the sick; to inquire after}

\label{lesson:TE-C519-paramarsincu}



\begin{quote}
Before the new one: say the Telugu for an anaesthetic, then the Telugu for a bandage.
\end{quote}

\subsection*{You'll want to know: \te{పరామర్శించు}}
\textbf{\te{పరామర్శించు}} — \emph{parāmarśiñcu} — \textquotedblleft{}to visit the sick; to inquire after\textquotedblright{}.

\begin{quote}
\te{ఆసుపత్రిలో} \te{ఉన్న} \te{స్నేహితుడిని} \te{మేము} \te{పరామర్శించాము}.
\end{quote}

\begin{quote}
\emph{āsupatrilō unna snēhituḍini mēmu parāmarśiñcāmu.}
\end{quote}

\begin{quote}
We visited our friend in hospital.
\end{quote}

Past (I …): \te{పరామర్శించాను}, \emph{parāmarśiñcānu}.

Also calling on a family in mourning to share their grief.

\begin{cousinweb}
\textbf{\te{పరామర్శించు}} (\emph{parāmarśiñcu}) is built on Sanskrit \emph{parāmarśa} (Monier-Williams), first \textquotedblleft{}a seizing, a touching\textquotedblright{}, with the Telugu verb ending -\te{ఇంచు}; in Telugu it means going to see someone who is ill or grieving.
\end{cousinweb}

\begin{grammarlens}[title={What you've built}]
One more word for describing what you lived and felt.
\end{grammarlens}

\subsection*{Your turn}
\begin{itemize}
\raggedright
  \item \emph{Say it:} \emph{parāmarśiñcu}
  \item \emph{Say it:} \emph{parāmarśiñcu}, once more
  \item \emph{Say it:} \emph{āsupatrilō unna snēhituḍini mēmu parāmarśiñcāmu.} — We visited our friend in hospital.
  \item \emph{Recall:} say the Telugu for an anaesthetic, then the Telugu for a bandage, then say \emph{parāmarśiñcu} again
\end{itemize}

\begin{tcolorbox}[breakable,colback=teal!4,colframe=teal!35!black,title={Before you move on}]
What does \te{పరామర్శించు} mean? (\textquotedblleft{}To visit the sick; to inquire after\textquotedblright{}.) Say it once more.
\end{tcolorbox}
